\documentclass[12pt,a4paper]{report}

% Packages
\usepackage[utf8]{inputenc}
\usepackage[T1]{fontenc}
\usepackage[french]{babel}
\usepackage{geometry}
\usepackage{graphicx}
\usepackage{fancyhdr}
\usepackage{titlesec}
\usepackage{xcolor}
\usepackage{listings}
\usepackage{hyperref}
\usepackage{array}
\usepackage{booktabs}
\usepackage{enumitem}
\usepackage{tocloft}
\usepackage{setspace}

% Marges
\geometry{top=2.5cm, bottom=2.5cm, left=3cm, right=2.5cm}

% Couleurs
\definecolor{bleuFonce}{RGB}{0, 51, 102}
\definecolor{bleuClair}{RGB}{0, 102, 204}
\definecolor{grisCode}{RGB}{245, 245, 245}
\definecolor{vertCode}{RGB}{0, 128, 0}
\definecolor{rougeCode}{RGB}{180, 0, 0}

% Style des titres
\titleformat{\chapter}[block]
  {\normalfont\Large\bfseries\color{bleuFonce}}
  {\thechapter.}{1em}{}
\titleformat{\section}
  {\normalfont\large\bfseries\color{bleuClair}}
  {\thesection}{1em}{}
\titleformat{\subsection}
  {\normalfont\normalsize\bfseries}
  {\thesubsection}{1em}{}

% Style du code
\lstset{
  backgroundcolor=\color{grisCode},
  basicstyle=\ttfamily\small,
  keywordstyle=\color{bleuClair}\bfseries,
  commentstyle=\color{vertCode},
  stringstyle=\color{rougeCode},
  breaklines=true,
  frame=single,
  rulecolor=\color{gray},
  numbers=left,
  numberstyle=\tiny\color{gray},
  language=Java,
  showstringspaces=false,
  tabsize=4
}

% En-tête et pied de page
\pagestyle{fancy}
\fancyhf{}
\fancyhead[L]{\small Jeu de l'Espace -- Space Defender}
\fancyhead[R]{\small Projet JavaFX}
\fancyfoot[C]{\thepage}
\renewcommand{\headrulewidth}{0.4pt}

% Espacement
\onehalfspacing

\begin{document}

% ===================== PAGE DE GARDE =====================
\begin{titlepage}
  \centering

  \vspace*{1cm}

  {\Large \textbf{Université -- Département Informatique}}\\[0.3cm]
  {\large Année universitaire 2024--2025}

  \vspace{1.5cm}
  \rule{\linewidth}{1.5pt}\\[0.4cm]
  {\Huge \textbf{\color{bleuFonce} Rapport de Projet}}\\[0.3cm]
  {\LARGE \textbf{Sujet 1 : Jeu de l'Espace}}\\[0.2cm]
  {\Large \textit{Space Defender -- Jeu en JavaFX}}\\[0.2cm]
  \rule{\linewidth}{1.5pt}

  \vspace{2cm}

  {\large \textbf{Réalisé par :}}\\[0.5cm]
  \begin{tabular}{c}
    \large Hamza Sadik \\[0.3cm]
    \large Jaffal Mohammed \\[0.3cm]
    \large Youssef Sabri \\
  \end{tabular}

  \vspace{2cm}

  {\large \textbf{Encadrant :} Pr. [Nom de l'encadrant]}

  \vspace{2cm}

  {\large \today}

\end{titlepage}

% ===================== TABLE DES MATIÈRES =====================
\tableofcontents
\newpage

% ===================== CHAPITRE 1 : INTRODUCTION =====================
\chapter{Introduction}

\section{Contexte général}
Dans le cadre du module de développement d'applications graphiques, nous avons été chargés de concevoir et développer un jeu vidéo interactif intitulé \textbf{Space Defender}. Ce projet s'inscrit dans le sujet numéro 1 proposé : \textit{Développement d'un jeu de l'espace en JavaFX}.

Ce jeu s'inspire des classiques du genre \textit{Space Invaders}, où le joueur doit défendre la Terre contre une invasion d'extraterrestres qui larguent des bombes destructrices.

\section{Objectifs du projet}
Les principaux objectifs de ce projet sont :
\begin{itemize}[leftmargin=1.5cm]
  \item Développer une application graphique interactive en \textbf{JavaFX}
  \item Mettre en \oe{}uvre les concepts d'animation et de boucle de jeu en temps réel
  \item Gérer les interactions clavier et les collisions entre objets
  \item Concevoir une interface ergonomique et immersive
  \item Implémenter un système de niveaux progressifs
\end{itemize}

\section{Organisation du rapport}
Ce rapport est structuré en quatre chapitres :
\begin{enumerate}[leftmargin=1.5cm]
  \item \textbf{Introduction} : contexte et objectifs
  \item \textbf{Analyse et conception} : règles du jeu, architecture et diagrammes
  \item \textbf{Implémentation} : description des classes et extraits de code
  \item \textbf{Conclusion} : bilan et perspectives
\end{enumerate}

% ===================== CHAPITRE 2 : CONCEPTION =====================
\chapter{Analyse et Conception}

\section{Présentation du jeu}
\textbf{Space Defender} est un jeu de tir spatial à défilement vertical. Le joueur contrôle une fusée en bas de l'écran et doit détruire les OVNIs ennemis qui avancent vers lui tout en évitant leurs bombes.

\subsection{Règles du jeu}
\begin{itemize}[leftmargin=1.5cm]
  \item Le joueur dispose de \textbf{3 vies} représentées par des c\oe{}urs
  \item Il perd une vie chaque fois qu'une bombe ennemie le touche
  \item Le jeu est \textbf{gagné} quand tous les OVNIs sont détruits
  \item Le jeu est \textbf{perdu} quand toutes les vies sont épuisées
  \item Chaque niveau augmente la vitesse et le nombre d'ennemis
\end{itemize}

\subsection{Contrôles du joueur}
\begin{center}
\begin{tabular}{|c|l|}
\hline
\textbf{Touche} & \textbf{Action} \\
\hline
$\leftarrow$ / Q & Déplacer la fusée à gauche \\
$\rightarrow$ / D & Déplacer la fusée à droite \\
ESPACE & Tirer un projectile \\
ENTRÉE & Rejouer / Niveau suivant \\
ÉCHAP & Retour au menu \\
\hline
\end{tabular}
\end{center}

\section{Architecture du projet}

\subsection{Structure des fichiers}
Le projet est organisé selon l'architecture Maven standard :

\begin{verbatim}
SpaceGame/
├── pom.xml
└── src/
    └── main/
        └── java/
            └── com/spacegame/
                ├── MainApp.java
                ├── MenuScreen.java
                ├── GameScreen.java
                ├── Player.java
                ├── Enemy.java
                ├── Bullet.java
                ├── Bomb.java
                ├── Explosion.java
                └── Star.java
\end{verbatim}

\subsection{Diagramme des classes}
Le projet contient \textbf{9 classes} principales organisées ainsi :

\begin{center}
\begin{tabular}{|l|l|}
\hline
\textbf{Classe} & \textbf{Rôle} \\
\hline
\texttt{MainApp} & Point d'entrée de l'application JavaFX \\
\texttt{MenuScreen} & Écran d'accueil avec boutons Jouer/Quitter \\
\texttt{GameScreen} & Boucle principale du jeu (logique + rendu) \\
\texttt{Player} & La fusée contrôlée par le joueur \\
\texttt{Enemy} & Les OVNIs ennemis \\
\texttt{Bullet} & Les projectiles tirés par le joueur \\
\texttt{Bomb} & Les bombes larguées par les ennemis \\
\texttt{Explosion} & Animation d'explosion à l'impact \\
\texttt{Star} & Les étoiles du fond animé \\
\hline
\end{tabular}
\end{center}

\section{Boucle de jeu}
Le jeu utilise un \texttt{AnimationTimer} de JavaFX pour exécuter la boucle de jeu à environ 60 images par seconde. Chaque frame effectue :

\begin{enumerate}[leftmargin=1.5cm]
  \item \textbf{Mise à jour} (\texttt{update}) : déplacement des entités, détection des collisions, vérification des conditions de fin
  \item \textbf{Rendu} (\texttt{render}) : dessin du fond, des entités, du HUD et des effets
\end{enumerate}

% ===================== CHAPITRE 3 : IMPLÉMENTATION =====================
\chapter{Implémentation}

\section{Technologies utilisées}
\begin{itemize}[leftmargin=1.5cm]
  \item \textbf{Java 17} : langage de programmation principal
  \item \textbf{JavaFX 21} : bibliothèque graphique pour l'interface et les animations
  \item \textbf{Maven} : outil de gestion des dépendances et de build
  \item \textbf{IntelliJ IDEA} : environnement de développement
\end{itemize}

\section{Classe principale -- MainApp}
La classe \texttt{MainApp} est le point d'entrée de l'application. Elle initialise la fenêtre et charge l'écran de menu.

\begin{lstlisting}[caption={MainApp.java -- Point d'entrée}]
public class MainApp extends Application {
    public static final int WIDTH = 900;
    public static final int HEIGHT = 600;

    @Override
    public void start(Stage stage) {
        MenuScreen menu = new MenuScreen(stage);
        Scene scene = new Scene(menu.getRoot(), WIDTH, HEIGHT);
        stage.setTitle("Space Defender");
        stage.setScene(scene);
        stage.show();
    }
}
\end{lstlisting}

\section{Gestion du joueur -- Player}
La classe \texttt{Player} représente la fusée du joueur. Elle gère le dessin, le mouvement et les vies.

\begin{lstlisting}[caption={Player.java -- Mouvement et vies}]
public void moveLeft(double speed) {
    x = Math.max(0, x - speed);
}

public void moveRight(double speed) {
    x = Math.min(MainApp.WIDTH - width, x + speed);
}

public void loseLife() { lives--; }
public boolean isDead() { return lives <= 0; }
\end{lstlisting}

\section{Comportement des ennemis -- Enemy}
Les OVNIs se déplacent horizontalement et rebondissent sur les bords. À chaque rebond, ils descendent de 20 pixels. Ils larguent des bombes périodiquement.

\begin{lstlisting}[caption={Enemy.java -- Déplacement et bombes}]
public void update(double delta) {
    x += speedX * delta;
    if (x <= 0 || x + width >= MainApp.WIDTH) {
        speedX = -speedX;
        y += 20;  // Descend a chaque rebond
    }
    bombTimer += delta;
}

public boolean shouldDropBomb() {
    if (bombTimer >= bombCooldown) {
        bombTimer = rand.nextDouble() * bombCooldown * 0.5;
        return true;
    }
    return false;
}
\end{lstlisting}

\section{Détection des collisions}
La détection de collision utilise la méthode des \textbf{boîtes englobantes} (\textit{AABB -- Axis-Aligned Bounding Box}) :

\begin{lstlisting}[caption={Méthode de détection de collision}]
public boolean intersects(double ox, double oy,
                           double ow, double oh) {
    return x < ox + ow && x + width > ox
        && y < oy + oh && y + height > oy;
}
\end{lstlisting}

Deux types de collisions sont gérés :
\begin{itemize}[leftmargin=1.5cm]
  \item \textbf{Projectile $\leftrightarrow$ Ennemi} : l'ennemi est détruit et une explosion apparaît
  \item \textbf{Bombe $\leftrightarrow$ Joueur} : le joueur perd une vie et une explosion apparaît
\end{itemize}

\section{Système de niveaux}
La difficulté augmente à chaque niveau :

\begin{center}
\begin{tabular}{|c|c|c|c|}
\hline
\textbf{Niveau} & \textbf{Lignes d'ennemis} & \textbf{Vitesse (px/s)} & \textbf{Délai bombes (s)} \\
\hline
1 & 3 & 90 & 2.5 \\
2 & 4 & 120 & 2.0 \\
3 & 5 & 150 & 1.5 \\
4+ & 6+ & 180+ & 1.0 \\
\hline
\end{tabular}
\end{center}

\section{Effets visuels}
\subsection{Fond étoilé animé}
80 étoiles de tailles et vitesses aléatoires défilent vers le bas pour simuler le déplacement dans l'espace.

\subsection{Explosions}
À chaque destruction, une animation d'explosion est créée avec un cercle orange/jaune qui grandit puis disparaît progressivement en 0.4 secondes.

\subsection{Dessin des entités}
Toutes les entités sont dessinées avec l'API \texttt{GraphicsContext} du \texttt{Canvas} JavaFX, sans utilisation d'images externes.

% ===================== CHAPITRE 4 : CONCLUSION =====================
\chapter{Conclusion}

\section{Bilan du projet}
Ce projet nous a permis de développer un jeu vidéo complet en Java avec JavaFX. Nous avons mis en \oe{}uvre les concepts suivants :

\begin{itemize}[leftmargin=1.5cm]
  \item La \textbf{boucle de jeu} avec \texttt{AnimationTimer} et gestion du temps (\texttt{delta})
  \item Le \textbf{dessin 2D} avec \texttt{Canvas} et \texttt{GraphicsContext}
  \item La \textbf{gestion des événements clavier} en temps réel
  \item La \textbf{détection de collisions} par boîtes englobantes (AABB)
  \item L'\textbf{architecture orientée objet} avec séparation des responsabilités
  \item Un \textbf{système de niveaux} progressifs
\end{itemize}

\section{Difficultés rencontrées}
Parmi les difficultés rencontrées lors du développement :
\begin{itemize}[leftmargin=1.5cm]
  \item La synchronisation de la boucle de jeu avec JavaFX
  \item La gestion correcte des événements clavier entre les scènes
  \item L'encodage des caractères spéciaux dans l'affichage canvas
\end{itemize}

\section{Perspectives d'amélioration}
Des améliorations futures pourraient inclure :
\begin{itemize}[leftmargin=1.5cm]
  \item Ajout d'effets sonores (tirs, explosions, musique de fond)
  \item Intégration d'images et de sprites pour les personnages
  \item Un système de sauvegarde des meilleurs scores
  \item Des power-ups (bouclier, tir multiple, etc.)
  \item Un mode multijoueur local
\end{itemize}

\section{Conclusion générale}
Ce projet nous a permis d'acquérir une expérience pratique solide dans le développement d'applications graphiques interactives avec JavaFX. La réalisation d'un jeu complet, de la conception à l'implémentation, nous a appris à gérer la complexité d'une application temps réel tout en appliquant les principes de la programmation orientée objet.

\end{document}
